\documentclass[11pt]{article}

\usepackage[utf8]{inputenc}
\usepackage[T1]{fontenc}
\usepackage{lmodern}
\usepackage[margin=1in]{geometry}
\usepackage{amsmath,amssymb}
\usepackage{microtype}
\usepackage[hidelinks]{hyperref}

\newcommand{\code}[1]{\texttt{#1}}
\renewcommand{\thesection}{S\arabic{section}}

\title{Supplementary Material\\
\large Design rationale for spatial pathway-activity analysis:\\
the error-versus-coupling tension}
\author{}
\date{}

\begin{document}
\maketitle

This note gives the design rationale behind the linear-space, scale-invariant-loss
factorisation summarised in the Methods of the main text. It motivates, from first
principles, why a faithful spatial pathway-activity decomposition must \emph{not}
log-transform the data, and what it should optimise instead.

\section{Why a faithful decomposition cannot log-transform}
\label{s:rationale}

We ask a question logically prior to enrichment: \textbf{what should the factorisation
optimise so that its components are faithful to how pathways and measurements actually
behave?} Two properties of spatial-omics data constrain the answer; together they rule out
the conventional ``log-transform, then PCA/NMF'' recipe.

\paragraph{(R1) Measurement error is multiplicative.} Imaging-mass-spectrometry ion
intensities are corrupted by gain/efficiency fluctuations that scale with signal: a measured
value behaves like $x = \mu\,\varepsilon$ with multiplicative $\varepsilon > 0$, so error
grows with abundance. Sequencing-count modalities share this signal-dependent character ---
sampling (Poisson, and typically over-dispersed) noise whose variance grows with the mean.
The conventional remedy is to take logarithms, after which the error is approximately additive
and ordinary least-squares / PCA apply.

\paragraph{(R2) Pathway coupling is multiplicative.} Membership in a common pathway means
coordinated, \emph{proportional} change: if two features A and B (metabolites or transcripts,
say) are co-regulated, doubling the local activity of the pathway scales both. In linear space
a rank-1 component $u\,v^{\!\top}$ encodes precisely this --- the entry $u_p v_m$ is a
per-location activity times a per-feature loading, a product. This is the structure we want a
component to capture.

\paragraph{The logarithm cannot serve R1 and R2 simultaneously.} Under a log transform the
shared activity $s$ of a rank-1 pathway ($A = a\,s$, $B = b\,s$) becomes an \emph{additive}
offset, $\log A = \log a + \log s$: the multiplicative coupling that defines the pathway is
converted into additive-shared-log structure, a weaker and different statement. A linear
factorisation of log-transformed data therefore studies the wrong object for R2, even as the
log was introduced to fix R1. Conversely, staying in linear space to keep R2 intact leaves
ordinary least squares mismatched to the multiplicative error of R1. The conventional pipeline
thus contains an internal contradiction: the log can serve the error model \emph{or} the
coupling model, but not both.

\paragraph{Resolution.} \code{colorpath} resolves this rather than splitting the difference.
It keeps the data in linear space --- so the outer product is genuine multiplicative coupling
(R2) --- and obtains the multiplicative-error property (R1) from a \textbf{scale-invariant
loss function} (the linear-space Itakura--Saito / Kullback--Leibler route described in the
main-text Methods) instead of from a transform, satisfying both requirements simultaneously.
On top of this fidelity term, an optional Hilbert--Schmidt independence penalty can drive the
components toward ICA-style statistical independence without leaving the non-negative,
multiplicative model. For spatial transcriptomics the same logic forbids the routine
$\log(1+x)$ normalisation: counts are kept linear (a linear library-size correction only) and
the KL/IS loss carries the error model, so the regional sub-programmes a pathway resolves into
remain genuine multiplicative outer products.

\end{document}
